% unidades/05-rutina-diaria.tex
\chapter{⏰ 毎日の生活 — Rutina Diaria}
\unidadheader{05}{毎日の生活}{Rutina Diaria}{Hablar de actividades cotidianas}

\begin{tcolorbox}[vocabbox, title={📌 Vocabulario Nuevo}]
\begin{tabularx}{\linewidth}{llX}
\toprule
\textbf{Japonés} & \textbf{Español} & \textbf{Tipo} \\
\midrule
\vocabitem{起きる}{おきる}{levantarse}{v. gr.2}
\vocabitem{寝る}{ねる}{dormir / acostarse}{v. gr.2}
\vocabitem{食べる}{たべる}{comer}{v. gr.2}
\vocabitem{飲む}{のむ}{beber}{v. gr.1}
\vocabitem{働く}{はたらく}{trabajar}{v. gr.1}
\vocabitem{勉強する}{べんきょうする}{estudiar}{v. irr.}
\vocabitem{帰る}{かえる}{regresar (a casa)}{v. gr.1}
\vocabitem{毎日}{まいにち}{todos los días}{adv.}
\vocabitem{毎朝}{まいあさ}{todas las mañanas}{adv.}
\vocabitem{毎晩}{まいばん}{todas las noches}{adv.}
\bottomrule
\end{tabularx}
\end{tcolorbox}

\begin{tcolorbox}[phrasebox, title={⚙️ Conjugación: Forma ます (repaso grupos)}]
\begin{tabular}{llll}
\toprule
\textbf{Grupo} & \textbf{Regla} & \textbf{Ejemplo} & \textbf{ます形} \\
\midrule
1 (う-verbos) & う→い + ます & \ruby{飲}{の}む & \ruby{飲}{の}みます \\
2 (る-verbos) & る → ます & \ruby{食}{た}べる & \ruby{食}{た}べます \\
3 (irregulares) & memorizar & する / くる & します / きます \\
\bottomrule
\end{tabular}
\end{tcolorbox}

\subsection*{🗣️ Frases Clave}

\frase{\ruby{毎朝}{まいあさ}\ruby{六時}{ろくじ}に\ruby{起}{お}きます。}
  {Me levanto a las 6 todas las mañanas.}
\frase{\ruby{学校}{がっこう}で\ruby{昼}{ひる}ごはんを\ruby{食}{た}べます。}
  {Como el almuerzo en la escuela.}
\frase{\ruby{夜}{よる}\ruby{十一時}{じゅういちじ}に\ruby{寝}{ね}ます。}
  {Me duermo a las 11 de la noche.}
\frase{\ruby{毎日}{まいにち}\ruby{日本語}{にほんご}を\ruby{勉強}{べんきょう}します。}
  {Estudio japonés todos los días.}

\subsection*{⚙️ Gramática}

\patron{Forma ます — afirmativa, negativa y pasada}
  {V-ます / V-ません / V-ました / V-ませんでした}
  {\ej{\ruby{毎日}{まいにち}\ruby{コーヒー}{こーひー} को \ruby{飲}{の}みます。}
    {Bebo café todos los días.}
  \ej{きのう\ruby{学校}{がっこう}に\ruby{行}{い}きませんでした。}
    {Ayer no fui a la escuela.}}

\patron{Hora de la acción}
  {〜時に〜ます}
  {\ej{\ruby{八時}{はちじ}に\ruby{起}{お}きます。}{Me levanto a las 8.}
  \ej{\ruby{七時}{しちじ}に\ruby{家}{いえ}に\ruby{帰}{かえ}ります。}
    {Regreso a casa a las 7.}}

\begin{tcolorbox}[ejerciciobox, title={📝 Ejercicios Rápidos}]
\begin{enumerate}
  \item Conjuga en ます形: \ruby{働}{はたら}く → ?
  \item Conjuga en negativa pasada: \ruby{食}{た}べる → ?
  \item Escribe tu rutina mañanera en 2 oraciones.
\end{enumerate}
\vspace{0.4em}
\textbf{Respuestas:}
1) \ruby{働}{はたら}きます\quad
2) \ruby{食}{た}べませんでした\quad
3) (personal)
\end{tcolorbox}

\begin{tcolorbox}[kanjibox, title={🀄 Kanjis de la Unidad}]
\begin{tabularx}{\linewidth}{cllXX}
\toprule
\textbf{Kanji} & \textbf{On} & \textbf{Kun} & \textbf{Significado} & \textbf{Vocab} \\
\midrule
\kanjirow{起}{き}{お}{levantarse}{\ruby{起}{お}きる / \ruby{起床}{きしょう}}
\kanjirow{寝}{しん}{ね}{dormir / acostarse}{\ruby{寝}{ね}る / \ruby{就寝}{しゅうしん}}
\kanjirow{食}{しょく}{た}{comer / alimento}{\ruby{食}{た}べる / \ruby{食事}{しょくじ}}
\kanjirow{飲}{いん}{の}{beber}{\ruby{飲}{の}む / \ruby{飲み物}{のみもの}}
\kanjirow{毎}{まい}{—}{cada / todo}{\ruby{毎日}{まいにち} / \ruby{毎週}{まいしゅう}}
\kanjirow{勉}{べん}{—}{esforzarse}{\ruby{勉強}{べんきょう}}
\bottomrule
\end{tabularx}
\end{tcolorbox}
